% Included section; theorem environments and \Z, \N, \Con, \Pos are supplied
% by the containing report.  All artifact paths below are relative to
% integer41/.  This file deliberately contains no preamble.
\section{A quartic in 41 unrestricted integer unknowns}

The main construction in this section is an explicitly specified polynomial
\(Q_{41}\in\Z[x_1,\ldots,x_{41}]\) of total degree four.  It reduces the
45 integer variables and 28 quadratic residuals of the inspected source
quartic to 41 integer variables and 26 residuals.  The first saving comes
from recovering the sign of the radix \emph{after} the Pell equations have
forced a power of two.  The other three come from using a larger scale to
exclude a negative common sign, so that two positivity representations can
be replaced by one representation of a positive product.

\begin{theorem}
For the four fixed integer coefficients specified below, ordinary
mathematics formalizable in ZFC proves
\[
 \exists\mathbf x\in\Z^{41}\;Q_{41}(\mathbf x)=0
 \quad\Longleftrightarrow\quad
 \text{there exists a finite ZFC refutation}.
\]
Consequently ZFC proves that the absence of integer zeros of \(Q_{41}\) is
equivalent to \(\Con(\mathrm{ZFC})\).  Exact elimination of polynomially
defined coordinates also gives an equivalent sextic in 36 integers and an
equivalent octic in 31 integers.
\end{theorem}

The statement uses the inherited native compiler, Pell lemmas, and
calibration of the literal blank K215 machine to syntactic ZFC
refutability.  Those interfaces are identified below.  This is not a new
Lean or Rocq kernel certificate and does not assert a global minimum.
All artifact paths in this section are relative to
\path{integer41/}.  The inspected ProveIt snapshot is
\path{beb0192cb8dc36f78c369c080f2adfbc1887a486}.

\subsection{The complete fixed coefficient tuple}

The coefficient ports are fixed numerals, not quantified variables.  Use
the \emph{whole} source34 tuple named
\path{mixed-parity-unrounded-1064921}, including its shifted high mask.
Put
\[
 d=1064921=121\cdot8801,\qquad
 \rho=2^{121},\qquad B=2^d,\qquad \beta=B-1.
\]
Here is a finite recipe that fixes the other three ports completely.  Let
\(x_{it}\) be the ordered west, north, east, and south labels of tile
\(i\), in the literal 943-entry list
\path{upstream/wang943_tiles.json}.  Let
\[
 (o_0,o_1,o_2,o_3)=(1020,1000,980,960),\qquad
 e_t(j)=o_t+2j,\qquad W_j=2^{10+2j},
\]
and define the finite set
\[
 \mathcal E=\{0,\ldots,942\}
 \cup\{e_t(j):0\leq t<4,\ 0\leq j<10\}
 \cup\{1063,1064,1095,1096,2195,3296\}.
\]
For \(0\leq i<943\), set
\[
 c_i=1+\sum_{t=0}^3\sum_{j=0}^9
       \operatorname{bit}_j(x_{it}+512)W_{10t+j}
       +\sum_{j=40}^{43}W_j.
\]
The remaining coefficients are
\(c_{e_t(j)}=W_{10t+j}\),
\((c_{1095},c_{1096},c_{2195},c_{3296})=(W_{40},W_{41},W_{42},W_{43})\),
and \(c_{1063}=c_{1064}=0\).  Finally put
\[
 \mu=1022+\sum_{j=0}^{43}W_j,\qquad
 Z_0=(2^{119}-1)\sum_{e=8792}^{8798}\rho^e
       +(2^{27}-1)\rho^{8799},
\]
\begin{align*}
 D_C&=\rho^{40}+\rho^{41}+\rho^{1100}+\rho^{2200}
       +\sum_{e\in\mathcal E}c_e
            (\rho^{4393-e}+\rho^{5495-e}),\\
 D_R&=\rho,\\
 M_C&=B-1-\sum_{\substack{e\in\mathcal E\\e\ne1}}\rho^e
             -2\rho^{1063},\\
 M_{F,n}&=4+\mu(\rho^{4393}+\rho^{5495})
       +\sum_{j=0}^9(\rho^{1000+2j}+\rho^{1021+2j})
       +\rho^{2195}+\rho^{3296}+Z_0,\\
 K_0&=D_C+BD_R,\qquad M_F=M_{F,n}+\beta.
\end{align*}
In particular, the port \(M_F\) is \emph{not} \(M_{F,n}\).  This source
binding includes the exact native geometry and completeness controller
proved in \path{upstream/COEFFICIENT_PROOF.md}; agreement of four coarse
inequalities alone would not establish the compiler interface.

The actual coefficients satisfy
\[
 \beta\geq17,\qquad 0<M_C<\beta,
 \qquad \beta<M_F<2\beta,\qquad K_0\geq0.
\]
They are supplied as gzip-compressed minimal unsigned little-endian byte
strings in \path{coefficients/beta.bin.gz}, \path{coefficients/K.bin.gz},
\path{coefficients/MC.bin.gz}, and
\path{coefficients/MF_source.bin.gz}, respectively.  Their bit lengths are
1064921, 1065043, 1064921, and 1064922; their total uncompressed size is
532479 bytes.  The builder reconstructs all four integers from the literal
tile list and matches every full raw SHA-256 digest in
\path{upstream/PAYLOAD_BINDING.json}.  The file
\path{SOURCE_PINS.json} records the exact source Git blobs.  No coefficient
depends on waiting for a machine to halt.

\subsection{All variables and all 26 quadratic residuals}

There are 32 arithmetic coordinates
\[
\begin{gathered}
 J,F,\alpha,z,f,h,r,v,w_4,\tau,\eta,\zeta,y,\gamma,o,j,\\
 N,E,P,c,a,X,D,C_2,L,F_2,V,g,I,B_{\mathrm{aux}},\Delta,u
\end{gathered}
\]
and nine further coordinates, grouped as
\(\boldsymbol\xi=(\xi_1,\xi_2,\xi_3)\),
\(\boldsymbol\upsilon=(\upsilon_1,\upsilon_2,\upsilon_3)\), and
\(\boldsymbol\omega=(\omega_1,\omega_2,\omega_3)\).
Every one of these 41 coordinates ranges over \(\Z\).
The names \(v,w_4\) are the JSON coordinates
\path{v_scale}, \path{w_scale}; the triples correspond, in order, to
\path{FA_pos1}, \path{FA_pos2}, \path{FA_pos3};
\path{FC_pos1}, \path{FC_pos2}, \path{FC_pos3}; and
\path{eta_zeta_pos1}, \path{eta_zeta_pos2}, \path{eta_zeta_pos3}.

Use the following \emph{affine abbreviations}, which introduce no
additional unknowns:
\[
 q=\beta J+1,\quad k=\eta+\zeta,\quad C=q-F-\alpha,
 \quad s=2v+1,\quad i=2r+1,\quad H=4a+3.
\]
The positive ternary form is
\[
 \Pos(x,y,z)=1+x^2+y^2+z(z+1).
\]
It is strictly positive on \(\Z^3\) and represents every positive
integer.  For surjectivity, apply the three-square theorem to \(4n-3\):
exactly one of the three squares is odd, so division by four gives the
displayed form after permuting the coordinates.

Define the residuals literally as follows:
\begin{align*}
 r_1&=L-q^2+1, &
 r_2&=N-q(L+1),\\
 r_3&=X-4qw_4, &
 r_4&=a-E-sN,\\
 r_5&=E(X+1)-Xa, &
 r_6&=c-\eta-k(a-E),\\
 r_7&=P-k-E(c-\eta), &
 r_8&=\tau^2-P(P-k)-1,\\
 r_9&=D-X-ac-\gamma H, &
 r_{10}&=\Delta-a^2-H,\\
 r_{11}&=C_2-c^2, &
 r_{12}&=D^2-\Delta C_2-1.
\end{align*}
The remaining rows are
\begin{align*}
 r_{13}&=k-1-hE-(q-F)(N-q)+CL-(M_C+qM_F)J,\\
 r_{14}&=(K_0+4w_4)C-z(q-1)-F,\\
 r_{15}&=F_2-f^2,\\
 r_{16}&=V+c-fo,\\
 r_{17}&=V-cj-hE+k-1,\\
 r_{18}&=g+V^2-y^2,\\
 r_{19}&=I-iC_2,\\
 r_{20}&=B_{\mathrm{aux}}-I^2,\\
 r_{21}&=B_{\mathrm{aux}}-\Delta F_2+\Delta,\\
 r_{22}&=y^2-B_{\mathrm{aux}}g-1,\\
 r_{23}&=X-2u^2,\\
 r_{24}&=F\alpha-\Pos(\boldsymbol\xi),\\
 r_{25}&=FC-\Pos(\boldsymbol\upsilon),\\
 r_{26}&=\eta\zeta-\Pos(\boldsymbol\omega).
\end{align*}
Here \(\Pos(\boldsymbol\xi)\) means
\(\Pos(\xi_1,\xi_2,\xi_3)\), and likewise for the other triples.
The required single equation is
\[
 \boxed{\quad Q_{41}=\sum_{\nu=1}^{26}r_\nu^2=0.\quad}
\]
Over integers this is exactly the simultaneous vanishing of all 26 rows.
After expanding the affine abbreviations, every row has degree at most
two.  Only \(r_8\) contains \(\tau\), through its term \(\tau^2\), so
\([\tau^4]Q_{41}=1\).  Thus the degree is exactly four for the actual
fixed coefficient tuple.  The machine-readable specification is
\path{quartic41.json}.

\subsection{The inherited Pell interfaces and the proof boundary}

For \(b\geq2\), define the integer Pell sequences by
\[
 \chi_b(n)+\psi_b(n)\sqrt{b^2-1}
       =(b+\sqrt{b^2-1})^n\qquad(n\geq0).
\]
We use the ordinary Pell classification and the following two inherited
rank results, with their size hypotheses stated explicitly.

\begin{lemma}[Ordinary strong rank and signed auxiliary step-down]
Let \(b\geq2\), \(\delta=b^2-1\),
\(c_0=\psi_b(p)>0\), and let \(i\) be a nonzero odd integer.
If
\[
 c_0>b\delta^2,\qquad (ic_0^2)^2=\delta(f^2-1),
\]
then \(|f|=\chi_b(m)\) for an integer \(m\geq1\) satisfying
\(p\mid m\), \(c_0\mid m\), and \(m\geq c_0\).
Suppose also that \(c_0>2p,2R\), \(R>0\), and, with
\(\varrho=|i|c_0^2\), there are integers \(V,y\) satisfying
\[
 \varrho^2V^2-(\varrho^2-1)y^2=1,\qquad
 V\equiv\pm c_0\pmod{|f|},\qquad V\equiv-R\pmod{c_0}.
\]
Then \(p=R\).
\end{lemma}

The ordinary strong-rank input is the relaxed-rank lemma cited in
\path{upstream/SIGNED_DOMAIN_29.md}, under ``Ordinary-strong systems and
product''.  Its use for either sign of the scale is recorded in
\path{upstream/ODD_SCALE_27_PROOF.md}, Section 2.  For clarity, the second
part follows from the signed auxiliary proof as follows.  Auxiliary Pell
classification gives an odd index \(\ell=2j_0+1\) and an odd Chebyshev
quotient \(Q_{j_0}\) with
\[
 Q_{j_0}(1-b^2)=(-1)^{j_0}\psi_b(\ell),\qquad
 Q_{j_0}(0)=(-1)^{j_0}\ell.
\]
Reducing modulo \(|f|\), squaring the first exposed congruence, and
applying the plus step-down lemma with \(0<2p<m\) gives
\(\ell\equiv\pm p\pmod m\).  Reducing modulo \(c_0\) gives
\(\ell\equiv\pm R\pmod{c_0}\).  As \(c_0\mid m\), this yields
\(R\equiv\pm p\pmod{c_0}\).  The strict bounds
\(0<R,p<c_0/2\) force equality.  In particular, the sign of an original
coordinate \(c=\pm c_0\) is immaterial here.  One must establish the
ordinary strong-rank conclusion before inferring
\(f^2\equiv1\pmod{c_0}\); simply dividing its norm by \(\delta\) is not
a valid substitute.

The proofs below also use the ratio estimates and main-root congruences
in \path{upstream/ODD_SCALE_27_PROOF.md}, Sections 3--4, and in
\path{upstream/SIGNED_DOMAIN_29.md}, under ``Signed gamma''.  We record
the substitutions and the inequalities that make those applications
valid; no finite search is being substituted for an unbounded Pell lemma.

\subsection{A sign-free packing bound and a positive global scale}

At any common zero set
\[
 \begin{gathered}
 Q=|q|,\quad Y=sq^3,\quad A=a+2,\quad
 G=q(q-F)-C,\\
 R=G(q^2-1)+(M_C+qM_F)J.
 \end{gathered}
\]
The first 13 rows give, by polynomial identities,
\[
 \begin{gathered}
 N=q^3,\quad E=XY,\quad a=(X+1)Y,\quad c=kY+\eta,\quad
 \Delta=A^2-1,\quad k=R+1+hE,\\
 \tau^2-(XY^2)(XY^2+1)k^2=1,\qquad
 D^2-\Delta c^2=1,\qquad D=X+ac+\gamma H.
 \end{gathered}
\]
For example, substituting \(a=E+Y\) into \(r_5=0\) gives
\(E=XY\), without division.  Substituting \(L=q^2-1\) and
\(N=q^3\) into \(r_{13}\) gives \(k-1-hE-R\).
Rows 15--22 give the further exact identities
\[
 (ic^2)^2=\Delta(f^2-1),\qquad
 (ic^2)^2(V^2-y^2)+y^2=1,\qquad
 V=of-c=jc-R.
\]

The two guards \(F\alpha>0\) and \(FC>0\) force
\(F,\alpha,C\) to have the same nonzero sign.  Their sum is \(q\),
so \(|q|\geq3\).  Since \(q=\beta J+1\), the case \(J=0\) is
impossible, and \(Q\geq\beta-1\geq16\).  This does not yet prove
\(q>0\).

If \(q=Q>0\), all three summands are positive and
\[
 Q+2\leq G\leq Q^2-Q-1,\qquad
 0<(M_C+qM_F)J<(2Q+1)(Q-1).
\]
If \(q=-Q<0\), put \(F_0=-F\), \(C_0=-C\).  Then
\[
 G=Q(Q-F_0)+C_0,\quad 2Q+1\leq G\leq Q^2-2,
 \quad 0<(M_C+qM_F)J<2Q(Q+1).
\]
These follow directly from the fixed mask bounds and
\(J=(q-1)/\beta\).  In the positive branch the resulting upper bound is
\(Q^4-Q^3\); in the negative branch it is
\(Q^4-Q^2+2Q+2\).  Consequently both branches satisfy
\[
 \boxed{\quad Q(Q^2-1)<R<Q^4-2,\qquad R>3Q+3.\quad}
\]
Neither \(\eta,\zeta\) nor transport has been used.

Since \(s\) is odd, \(Y\ne0\) and \(|Y|\geq Q^3\).  Row 23 gives
\(X\geq0\).  If \(X=0\), then \(a=Y\) and
\(D=ac+\gamma H\).  Reducing the main norm modulo \(H\), using
\(\Delta=a^2+H\), gives \(0\equiv1\pmod H\), impossible because
\(|H|=|4Y+3|>1\).  Thus \(X>0\).  The new scale row now gives
\[
 X=4qw_4\geq4Q,\qquad
 |E|=X|Y|\geq4Q^4>3R+1.
\]
The extra factor four costs no variable: \(w_4\) replaces the old
coordinate \(w\) through the affine substitution \(w=4w_4\).

\subsection{Recovering positivity of eta and zeta}

The last guard gives \(\eta\zeta>0\), so their signs agree and neither
is zero.  Suppose for contradiction that both are negative.  Write
\[
\begin{gathered}
 \eta_0=-\eta>0,\quad\zeta_0=-\zeta>0,\quad
 K=\eta_0+\zeta_0=-k>0,\\
 Z=|Y|,\quad
 b=|A|,\quad c_0=|c|,\quad P_0=2XY^2+1.
\end{gathered}
\]
For \(Y>0\) one has \(KY<c_0<K(Y+1)\); for \(Y<0\) one has
\(K(Z-1)<c_0<KZ\).  In either case \(c_0>K\), \(b\geq2\), and
\(P_0>b\).  The last inequality follows, for example, from
\[
 2XZ^2+1-\big((X+1)Z+2\big)>0
 \quad(X\geq64,\ Z\geq4096).
\]
The first and main Pell classifications therefore give
\[
 K=2\psi_{P_0}(n),\quad |\tau|=\chi_{P_0}(n),\qquad
 c_0=\psi_b(p),\quad |D|=\chi_b(p),\qquad n,p\geq1.
\]
Monotonicity implies \(p>n\): if \(p\leq n\), then
\(c_0\leq\psi_b(n)\leq\psi_{P_0}(n)=K/2\), a contradiction.
Because \(P_0\equiv1\pmod{|E|}\), the Pell recurrence gives
\(\psi_{P_0}(n)\equiv n\pmod{|E|}\).  The index row and \(k=-K\)
give
\[
 2n\equiv-(R+1)\pmod{|E|}.
\]
As \(2n>0\) and \(0<R+1<|E|\), this entails
\[
 2n\geq |E|-R-1>2R,\qquad\text{hence }p>n>R.
\]
In particular \(p\geq6\).  The elementary Pell growth bound
\(\psi_b(p)\geq(2b-1)^{p-1}\) yields
\[
 c_0>b(b^2-1)^2,\qquad c_0>2p,\qquad c_0>2R.
\]
For the first inequality it suffices to observe that
\((2b-1)^5>b^5>b(b^2-1)^2\) for \(b\geq2\); the second follows
from \(3^{p-1}>2p\) for \(p\geq6\).

Every hypothesis of the ordinary strong-rank and auxiliary lemma now
holds.  In particular \(i=2r+1\) is nonzero and odd, the two exposed
quotients give \(V\equiv-c\pmod{|f|}\) and
\(V\equiv-R\pmod{c_0}\), and \(c=\pm c_0\) is permitted.  The
lemma gives \(p=R\), contradicting \(p>n>R\).  Therefore
\[
 \eta>0,\qquad\zeta>0.
\]
This argument is the reason for the larger scale.  The weaker bound
\(|E|>R+2\) from the unscaled construction does not supply the
negative-index contradiction \(n>R\).

\subsection{Delayed radix sign recovery: the intermediate quartic44}

It remains to justify that the other signed coordinates cannot introduce
spurious roots.  For this purpose consider the intermediate system
\path{quartic44.json}: replace \(w_4\) by one coordinate \(w\), use
\(X=qw\), and replace the product guard on \(\eta\zeta\) by two
separate ternary guards \(\eta=\Pos(\mathbf e)\),
\(\zeta=\Pos(\mathbf t)\).  It has 44 variables and 27 residuals.
All other rows are unchanged, in particular \(X=2u^2\).  Starting with
a zero of the 41-variable system, put \(w=4w_4\) and supply ternary
representations of the now positive \(\eta,\zeta\).  This gives a zero
of the intermediate system exactly.

For a zero of that intermediate system the preceding sign-free packing
bounds remain valid.  The same zero-scale exclusion gives
\[
 X\geq Q\geq16,\qquad |Y|\geq Q^3,\qquad
 |E|\geq Q^4>R+2,\qquad \eta,\zeta>0.
\]
These are the early scale hypotheses required by the inherited signed
Pell argument.  In particular no assumption that \(q>0\) or that
\(w\) is a square is needed to establish them.

Here are the rank and sign steps explicitly.  If \(Y<0\), put
\(Z=-Y\), \(b=(X+1)Z-2\), \(c_0=-c=kZ-\eta\), and
\(P_0=2XZ^2+1\).  Then
\[
 k(Z-1)<c_0<kZ,\quad P_0>b,\quad c_0>k,\qquad
 k=2\psi_{P_0}(n),\quad c_0=\psi_b(p),\quad p>n.
\]
Now the index congruence has the positive sign:
\(2n=R+1+t_0|E|\), \(t_0\geq0\).  Thus
\(n\geq(R+1)/2\geq24\).  Pell growth gives
\(c_0>b(b^2-1)^2\) and \(c_0>2p\); also
\(b=|E|+Z-2>R\), so \(c_0>2R\).
The inherited rank lemma and the exposed \(V\)-congruences give
\(p=R\).  Since \(n<p=R\) and \(|E|>R+2\), there is no wrap:
\(t_0=0\), hence \(R=2n-1\).

Writing \(R=2r_0+1\), the negative-scale ratio estimate gives
\[
 Z>\frac{c_0}{k}\geq
 \frac12\left(\frac{(2b-1)^2}{2P_0}\right)^{r_0}
 >\frac{X^{r_0}}2>2^R.
\]
The strict middle inequality follows from
\[
 (2b-1)^2-2P_0X
 =4(2X+1)Z^2-20(X+1)Z+25-2X>0;
\]
for \(X\geq16\) this is positive at \(Z=4\), where it equals
\(46X+9\), and increases thereafter.  Here \(r_0\geq2\).
The two signs of the main root imply, respectively,
\[
 X\equiv-2^R\pmod H
 \quad\hbox{or}\quad X2^R\equiv1\pmod H.
\]
Both are impossible, because
\(0<X+2^R<X+Z<|H|\) and \(1<X2^R<XZ<|H|\).
Consequently \(Y>0\).

For \(Y>0\), the same rank argument applies with \(b=A\) and
\(kY<c<k(Y+1)\).  In particular the index row implies \(h\geq0\),
the Pell indices satisfy \(p>n\geq(R+1)/2\), and the growth conditions
of the ordinary strong-rank lemma again hold.  Thus \(p=R=2n-1\).
The positive-scale ratio estimate gives
\[
 a>\frac{X^{r_0+1}}3,\qquad 0<X2^R<a<H-1.
\]
If \(D=-\chi_A(R)\), the main-root recurrence would give
\(X2^R\equiv-1\pmod H\), contradicting that interval.
Hence \(D=\chi_A(R)\), and the recurrence gives
\(X\equiv2^R\pmod H\).  Both \(X\) and \(2^R\) lie strictly
between zero and \(H\), so
\[
 \boxed{X=2^R.}
\]
This is precisely where the sign of the radix can finally be recovered.

\begin{lemma}[Mersenne radix sign]
For \(d\geq3\) and \(t\geq0\),
\(-2^t\not\equiv1\pmod{2^d-1}\).
\end{lemma}
\begin{proof}
Reduce \(t\) modulo \(d\), obtaining \(0\leq t_0<d\).
The least positive residue of \(2^t\) is \(2^{t_0}\), whereas
\(1\leq2^{t_0}\leq2^{d-1}<2^d-2\).  It therefore cannot be the
residue \(-1\).
\end{proof}

Since \(q\mid X=2^R\), one has \(q=\pm2^t\) for some
\(0\leq t\leq R\).  The definition \(q=\beta J+1\) and the lemma
exclude the negative sign.  Thus \(q>0\), and common signs give
\(F,\alpha,C>0\).  Also \(J>0\), \(w=X/q>0\), and
\(s=Y/q^3>0\).  This argument delays the radix sign until it is justified;
it does not attempt to recover the omitted square condition on \(w\).

\subsection{Literal transport, compiler soundness, and completeness}

The literal transport identity is
\[
 (K_0+w)C=z(q-1)+F,\qquad w=4w_4
\]
in the 41-variable system.  Source57's accompanying prose contains a
coordinate typo: it writes \((z-1)(q-1)+F\) at this point, but that is
not the expansion of its literal row.  Accordingly the conclusion needed
here is \(z>0\).  Indeed \(R>3q\) implies \(2^R>q^2\), so
\(w=2^R/q>q>F\).  With \(C\geq1\) and \(K_0\geq0\), transport
gives \(z(q-1)>0\).  In the older unshifted positive-factor convention,
the correct coordinate is \(z_{\mathrm{source}}=z+1\); its transport
factor then equals one.

After power recovery, the compiler soundness used here is the stage
``Power and compiler recovery'' and the blank finite-window stage of
\path{upstream/FIXED_BLANK_67_PROOF.md}, transported to the full tuple
above by \path{upstream/COEFFICIENT_PROOF.md}.  Its inputs are the
recovered positive coordinates, the authentic Pell kernel,
\(q^3\mid Y\), the bounds \(3q+1<R<q^4\), and the literal transport
identity.  The half-binomial valuation gives the population threshold;
\(q=\beta J+1\) and the fact that \(q\) is a power of two imply
\(q=B^{n_{\mathrm{cell}}}\) for some positive integer
\(n_{\mathrm{cell}}\); inverse population and
transport give the disjoint native masks.  The tuple's complete native
geometry decodes a blank K215 accepting tableau.  Its inherited finitary
calibration gives a finite ZFC refutation.

The square condition on \(w\) and the condition that \(q\) be twice a
square are sufficient choices of canonical witnesses in source57, not
hypotheses of this post-power soundness stage.  A current solution need
not satisfy those two restrictions in its existing coordinates.  To
recover a source57 solution one may first decode the finite computation
and then choose fresh canonical witnesses.  This distinction is necessary
for the delayed-sign reduction.

Conversely, a finite ZFC refutation supplies a blank K215 halt and the
canonical witnesses of the full source34 adapter.  They satisfy
\[
 q=2^t\ (t\text{ odd}),\quad R\equiv3\pmod4,\quad R>t,
 \quad X=2^R,\quad w=2^{R-t},
\]
with positive \(F,\alpha,C,\eta,\zeta\) and positive odd \(s,i\).
The adapter proves these parity promises for its actual coefficients,
including its rank controller; an older power-of-five assumption is not
being transferred to the new exponent \(d\).
Because \(R-t\) is a positive even integer, it is at least two, and
\(4\mid w\).  Choose
\[
 w_4=2^{R-t-2},\qquad u=2^{(R-1)/2},\qquad
 v=(s-1)/2,\qquad r=(i-1)/2.
\]
Retain the other canonical arithmetic coordinates and supply ternary
representations of \(F\alpha\), \(FC\), and \(\eta\zeta\).
Every displayed residual vanishes.  This proves completeness and, with
the preceding soundness argument, the theorem.  No consistency premise
is needed in either the fixed coefficient recipe or the equivalence.

\subsection{Exact sextic and octic companions}

For a sextic, delete the five indicated definition rows and eliminate
their five coordinates everywhere by
\[
 N=q(L+1),\quad X=4qw_4,\quad C_2=c^2,
 \quad F_2=f^2,\quad B_{\mathrm{aux}}=I^2.
\]
The remaining 21 residuals in 36 integer variables have degree at most
three.  Their sum of squares, specified in \path{sextic36.json}, has
exact degree six, with
\[
 [I^4g^2]Q_{36}=1.
\]
The certificate comes from the retained residual \(y^2-I^2g-1\);
the complete sparse expansion checks the coefficient in the full sum.

For an octic, instead delete ten definition rows and eliminate
\[
 \begin{gathered}
 L=q^2-1,\quad N=q^3,\quad X=4qw_4,\quad
 C_2=c^2,\quad F_2=f^2,\quad B_{\mathrm{aux}}=I^2,\\
 g=y^2-V^2,\qquad P=E(c-\eta)+\eta+\zeta,\\
 D=4qw_4+ac+\gamma(4a+3),\qquad \Delta=a^2+4a+3.
 \end{gathered}
\]
All right-hand sides involve only retained variables.  The remaining 16
residuals in 31 integers have degree at most four.  Their sum of squares
in \path{octic31.json} has exact degree eight, with
\[
 [I^4V^4]Q_{31}=1,
\]
as verified in the full expansion; the contributing auxiliary residual
is \(I^2(V^2-y^2)+y^2-1\).

These are exact bijections of integer zero sets: projection forgets the
defined coordinates, and their polynomial definitions restore them
uniquely.  Each deleted row is a unit-coefficient definition and becomes
the zero polynomial after substitution.  For \(N\), the simultaneous
octic substitution uses \(q(L+1)=q^3\).  Thus neither companion needs
a new sign, Pell, or compiler argument.

\subsection{Reproduction, comparison, and limitations}

Run \path{build_verify44.py}, then \path{build_verify41.py}, then
\path{build_companions.py} with Python 3; all three also run under
\texttt{python3 -O}.  They use the standard library, reconstruct the fixed
integers, check their full hashes and source blob identities, check the
packing and literal transport identities, and verify exact substitution
and elimination.  They expand the sums of squares with symbolic fixed
ports, check that every listed coordinate is live, and certify the exact
degrees and displayed unit coefficients.  The resulting certificates
are in \path{certificates/}.  Finite diagnostics for the elementary bounds
and radix lemma supplement their proofs above; they do not certify the
universal Pell or compiler theorems by sampling.

The precise primary comparison is the inspected source57 quartic:
45 variables and 28 rows become 41 variables and 26 rows at degree four.
The sextic and octic improve the named historical 42-variable sextic and
36-variable octic comparisons to 36 and 31 variables, respectively.
These claims do not displace the collection's smaller-variable equations
of much higher degree, and no global minimality or established worldwide
novelty is asserted.  The contribution is the explicit delayed-sign
guard, the absolute-radix bootstrap, and the large-scale exclusion of a
negative common sign, composed with the inherited full coefficient and
compiler interfaces.  The mathematical proof has internal AI-assisted
checking; no new end-to-end proof-assistant certificate or external peer
review is claimed.  Attribution and the exact inherited licenses are
preserved in \path{NOTICE.md} and \path{licenses/}.
